\documentclass[11pt]{scrartcl}


\usepackage[sexy]{evan}
\usepackage{import}
\usepackage[table]{xcolor}
\usepackage{xfp}

\newcommand{\gad}{\textcolor{yellow}{$\bigstar$}}
\newcommand{\bad}{\textcolor{red}{$\bigstar$}}

\definecolor{dcol0}{HTML}{C8E6C9}
\definecolor{dcol1}{HTML}{D4E9B3}
\definecolor{dcol2}{HTML}{E5ED9A}
\definecolor{dcol3}{HTML}{FFF59D}
\definecolor{dcol4}{HTML}{FFE082}
\definecolor{dcol5}{HTML}{FFCC80}
\definecolor{dcol6}{HTML}{FFAB91}
\definecolor{dcol7}{HTML}{F49890}
\definecolor{dcol8}{HTML}{E57373}
\definecolor{dcol9}{HTML}{D32F2F}

\makeatletter
\newcommand{\getcolorname}[1]{dcol#1}
\makeatother

\newcommand{\dif}[1]{%
    \edef\colorindex{\number\fpeval{floor(#1)}}%
    \edef\fulltext{#1}%
    \colorbox{\getcolorname{\colorindex}}{%
        \ifnum\colorindex>8
            \textbf{\textcolor{white}{\,\fulltext\,}}%
        \else
            \textbf{\textcolor{black}{\,\fulltext\,}}%
        \fi
    }%
}
% Variable para dificultad (inicial 0)
\newcommand{\thmdifficulty}{0}

% Comando para asignar dificultad antes del problema
\newcommand{\problemdiff}[1]{\renewcommand{\thmdifficulty}{#1}}

% Estilo del problema que incluye dificultad antes del título
\declaretheoremstyle[
    headfont=\color{blue!40!black}\normalfont\bfseries,
    headformat={%
      \dif{\thmdifficulty}\quad \NAME~\NUMBER\ifx\relax\EMPTY\relax\else\ \NOTE\fi
    },
    postheadspace=1em,
    spaceabove=8pt,
    spacebelow=8pt,
    bodyfont=\normalfont
]{problemstyle}

    \declaretheorem[style=problemstyle,name=Problema,sibling=theorem]{problema}
    \declaretheorem[style=problemstyle,name=Problema,numbered=no]{problema*}

\definecolor{yellow4}{RGB}{255, 206, 0}

\title {Orden}
\subtitle{Entrenamientos Nacionales Platas y EGMOs}

\author{Emmanuel Buenrostro}


\begin{document}
\maketitle

\section{Teoria}
En el entrenamiento anterior habiamos visto que 
\begin{theorem}
    [Teorema de Euler]
    Para $a,n$ enteros coprimos se cumple que 
    \[a ^{\varphi (n)} \equiv 1 \pmod n \]
\end{theorem}

Pero existen muchos otros enteros $k$ tales que $a^k \equiv 1 \pmod n$, entre ellos por ejemplo todos los multiplos de $\varphi(n)$. \\
Entonces nos vamos a centrar en lo siguiente:
\begin{definition}
   El \textit{orden} $\ord_n (a)$ es el menor entero positivo tal que $a^{\ord_n(a)} \equiv 1 \pmod n$
\end{definition}

Por ejemplo el orden de $2 \pmod 7$ es $3$ ya que $2^3 = 8 \equiv 1 \pmod 7$ y los anteriores no son 1. \\

Pero, ¿de qué nos sirve el orden?, pues principalmente para tener una divisibilidad: 
\begin{theorem}
    Si $x^k \equiv 1 \pmod n$ si y solo si $\ord_n (x) \mid k$.
\end{theorem}
\begin{proof}
    Digamos que $k = q \ord_n(x) +r$ con $0 \leq r < \ord_n(x)$ entonces 
    \[x^k = x^{q\ord_n(x)+r}= (x^{ord_n(x)})^k \cdot x^r \equiv x^r \pmod n\] 
    Entonces $x^r  \equiv 1 \pmod n$, pero si $r \neq 0$ entonces $0 < r < \ord_n(x)$ pero como $\ord_n(x)$ es el entero positivo minimo que cumple $x^{\ord_n(x)}\equiv 1 \pmod n$, entonces no se puede. \\
    Por lo que $r=0$ y $\ord_n(x) \mid k$. \\
    Y todos los multiplos de $\ord_n(x)$ cumplen porque si $k=q\ord_{n}(x)$ entonces 
    \[x^{k}=(x^{\ord_n(x)})^q \equiv 1^q \equiv 1 \pmod n\]
\end{proof}

En particular
\begin{theorem}
    $\ord_n(x) \mid \varphi(n)$ para todo $n,x$.
\end{theorem}
Esto es porque $x^{\varphi(n)}\equiv 1 \pmod n$ por Euler. \\

\begin{example}
    [Fermat Christmas Theorem]
    Si $p \mid x^2+y^2$ para $x,y$ enteros, entonces se cumple alguna de estas: 
    \begin{itemize}
        \item $p=2$
        \item $p \mid \gcd (x,y)$
        \item $p \equiv 1 \pmod 4$
    \end{itemize}
\end{example}

\begin{soln}
Si $p \neq 2$ y $p \nmid \gcd (x,y)$ entonces 
\begin{align*}
    x^2 &\equiv -y^2 \pmod p \\
    \left(\frac xy\right)^2 &\equiv -1 \pmod p \\ 
    \left(\frac xy\right)^4 &\equiv 1 \pmod p 
\end{align*}
Entonces si $d= \ord_p(\frac xy)$ sabemos que $d \mid 4$ y $d\nmid 2$ por lo que $d = 4$ .\\
Tambien sabemos que $d \mid \varphi(p)= p-1$ entonces $p \equiv 1 \pmod 4$.
\end{soln}

\begin{example}
Encuentra todos los enteros positivos $n$ tales que $n \mid 2^n-1$.
\end{example}

\begin{soln}
Si $n>1$ sea $p$ el primo más pequeño que divide a $n$.  \\
Entonces 
\[2^n \equiv 1 \pmod p \]
Y 
\begin{align*}
    \ord_p (2) &\mid n \\
    \ord_p (2) &\mid p-1 
\end{align*}
Entonces 
\[\ord_p(2) \mid \gcd (n, p-1)\]
Y si $d= \gcd (n, p-1)$ y $d>1$ entonces $d$ es un divisor de $n$ que es al menos $p$, ya que este es el factor primo más chico. \\
Pero $ p \leq d \mid p-1 \Rightarrow p \leq p-1$ lo cual es falso, entonces $d=1$, y $p \mid 1$, por lo que no existe tal $p$ y $n=1$.
\end{soln}

\begin{theorem}
    Si $n \mid a^k+1$ entonces $\nu_2(2k)=\nu_2(\ord_n(a))$.
\end{theorem}
\begin{proof}
Esto es debido a que 
\begin{align*}
    a^k &\equiv -1 \pmod n \\
    a^{2k} &\equiv 1 \pmod n
\end{align*}
Entonces 
\begin{align*}
    \ord_n(a) &\mid 2k \\
    \ord_n(a) &\nmid k
\end{align*}
La parte impar de $\ord_n(a)$ como divide a $2k$ tambien divide a $k$, y entonces si $\nu_2(\ord_n(a)) \leq \nu_2(k)$ tendriamos que $\ord_n(a) \mid k$, por lo que $\nu_2(\ord_n(a))= \nu_2(2k)$.
\end{proof}


\begin{example}
   Demuestra que si $n^2 \mid 2^n+1$ entonces $3\mid n$.
\end{example}

\begin{soln}

Sea $p$ el divisor primo más pequeño de $n$, como $2^n+1$ es impar entonces $n$ también y $p\geq 3$, entonces 
\begin{align*}
2^n &\equiv -1 \pmod p^2 \\
2^{2n} &\equiv 1 \pmod p^2 
\end{align*}
Y $\nu_2(\ord_{p^2}(2))=\nu_2(2n)=1$. \\

Entonces 
\begin{align*}
    \ord_{p^2}(2) &\mid 2n \\ 
    \ord_{p^2}(2) &\mid \varphi(p^2)=p(p-1)
\end{align*}
Entonces
\[\ord_{p^2}(2) \mid \gcd (2n, p(p-1))\]
Como $\gcd (p-1, n)=1$ entonces 
\[\ord_{p^2}(2) \mid 2p\]
Y $\ord_{p^2}(2) =2, 2p$. \\
Si $ord_{p^2}(2)=2$ entonces $p^2 \mid 2^2-1=3$, contradicción. \\
Entonces $\ord_{p^2}(2)=2p$ y 
\begin{align*}
    

2^{2p} &\equiv 1 \pmod {2p} \\
4 \equiv 2^{2p} &\equiv 1 \pmod {p} 
\end{align*}
Entonces $p \mid 3$ y $p=3$.
\end{soln}

\newpage 
\section{Problemas}
\epigraph{creo que si vas a ser el entrenador con mas sesiones \\ :O}{Lalo} 
\subsection{Nivel I}
\problemdiff{3}
\begin{problem}
    Prueba que todos los factores primos de $2^p-1$ son mas grandes que $p$
\end{problem}

\problemdiff{3}
\begin{problem}
    Prueba que cualquier factor de $2^{2^{n}}+1$ es $1 \pmod{ 2^{n+1}}$
\end{problem}

\problemdiff{3}
\begin{problem}
    Sean $a,b$ enteros coprimos entre si y $n$ un entero positivo. Prueba que cualquier divisor impar de $a^{2^n}+b^{2^n}$ es congruente a $1 \pmod 2^{n+1}$.
\end{problem}

\problemdiff{3}
\begin{problem}
    Una sucesión $x_1,x_2,\ldots$ de enteros cumple que $x_1 \in \{5,7\}$ y $x_{k+1}\in \{5^{x_k}, 7^{x_k}\}$ para cada $k \geq 1$. ¿Cuáles son los posibles residuos de $x_{2012}$ al dividirlo entre 100.
\end{problem}

\subsection{Nivel II}
\problemdiff{4}
\begin{problem}
    %[IZHO 2020/1]
    Sea $n$ un entero positivo tales que para cualesquiera enteros positivos $a,b$ el número $2^{a}3^b+1$ no es divisible entre $n$. Prueba que $2^c+3^d$ tampoco es divisible entre $n$ para cualesquiera enteros positivos $c$ y $d$.
\end{problem}

\problemdiff{4}
\begin{problem}
    Sea $a\geq 2, n \geq 1$ enteros. Prueba que $n$ divide a $\varphi(a^n-1)$.
\end{problem}


\problemdiff{4}
\begin{problem}
    %[HMMT 2016 T2]
    Para cada entero positivo $n$, sea $c_n$ el entero positivo más pequeño tal que $n^{c_n}-1$ es divisible entre $210$, si ese entero positivo existe, si no $c_n=0$. ¿Cuánto vale $c_1+c_2+\ldots+c_{210}$?.
\end{problem}


\problemdiff{4}
\begin{problem}
    Demuestra que el número $3^n-2^n$ no es divisible entre $n$ para $n \geq 2$.
\end{problem}

\problemdiff{4.5}
\begin{problem}
    %[Bulgaria 1997]
    Determina todos los enteros positivos $m,n \geq 2$ tales que 
    \[\frac{1+m^{3^n}+m^{2\cdot 3^{n}}}{n}\]
    es un número entero.
\end{problem}

\subsection{Nivel III}

\problemdiff{5}
\begin{problem}
    Prueba que no hay raices primitivas modulo $2^n$ para $n \geq 3$. Esto es, demuestra que no existe un entero $g$ tal que $g,g^2,g^3, \ldots$ contiene a todos los residuos impares modulo $2^n$.
\end{problem}


\problemdiff{5}
\begin{problem}
    %[MOP 2011]
    Sea $p$ un primo y $n$ un entero positivo. Se cumple que $\nu_p (2^n-1)=1$. ¿Se debe cumplir que $\nu_p(2^{p-1}-1)=1$?
\end{problem}


\problemdiff{5}
\begin{problem}
    %[Schinzel]
    Encuentra todos los enteros $n$ tales que $n \mid 2^{n-1}+1$.
\end{problem}

\problemdiff{5.5}
\begin{problem}
    %[USA TST 2003/3]
    Encuentra todas las tercias de primos $(p,q,r)$ tal que $p\mid q^r+1, q\mid r^p+1$ y $r\mid p^q+1$.
\end{problem}

\subsection{Nivel IV}
\problemdiff{6}
\begin{problem}
    %[USA TST 2008/4]
    Prueba que $n^7+7$ no es un cuadrado perfecto para todo entero $n$.
\end{problem}

\problemdiff{6}
\begin{problem}
    %[Romania TST 1996]
    Encuentra todas las parejas de primos $(p,q)$ tales que $3pq \mid n^{3pq}-n$ para todo entero positivo $n$. 
\end{problem}

\problemdiff{6}
\begin{problem}
    %[China 2009/2]
    Encuentra todos los números primos $p,q$ tales que $pq \mid 5^p+5^q$.
\end{problem}

\problemdiff{7}
\begin{problem}
    %[USA EGMO TST 2019/3]
    Sea $n$ un entero positivo tal que el número 
    \[\frac{1^k+2^k+\ldots+n^k}{n}\]
    es un entero para todo $k\in \{2,3,\ldots, 99\}$. Prueba que $n$ no tiene divisores entre $2$ y $100$ inclusive. 
\end{problem}

\problemdiff{7}
\begin{problem}
    %[China TST 2005/18]
    Prueba que para $n>2$, el mayor factor primo de $2^{2^n}+1$ es mayor o igual a $n\cdot 2^{n+2}+1$
\end{problem}


\end{document}